// Optical Flow — Lucas-Kanade motion estimation
//
// Computes per-pixel motion vectors between two consecutive frames using
// the Lucas-Kanade method: a local least-squares solve of the optical
// flow constraint equation (Ix*u + Iy*v + It = 0) over a window.
//
// For each pixel, the algorithm:
//   1. Computes spatial gradients (Ix, Iy) from the current frame
//   2. Computes the temporal gradient (It) between current and previous frame
//   3. Accumulates the 2x2 structure tensor A^T*A and vector A^T*b
//      over a local window
//   4. Solves the 2x2 linear system for (u, v) displacement
//
// The structure tensor's eigenvalues indicate estimation quality:
// corners give strong flow, edges give ambiguous flow (aperture problem),
// flat regions give no flow. A confidence threshold suppresses noise
// in textureless areas.
//
// Output: vec4 where R = horizontal motion, G = vertical motion (in pixels),
//         B = confidence (0-1), A = 1.
// Motion is encoded as signed pixel displacement (not 0.5-centered).
// Connect @current and @previous as consecutive frames.

i$window_radius = 4;    // window half-size (4 = 9x9 window, 81 equations)
f$confidence_threshold = 0.001;  // minimum eigenvalue to trust the estimate

// ── Accumulate structure tensor over the local window ────────────
// A^T A = [[sum(Ix^2), sum(Ix*Iy)],
//          [sum(Ix*Iy), sum(Iy^2)]]
// A^T b = [-sum(Ix*It), -sum(Iy*It)]
float sIx2 = 0.0;   // sum of Ix^2
float sIy2 = 0.0;   // sum of Iy^2
float sIxIy = 0.0;  // sum of Ix*Iy
float sIxIt = 0.0;  // sum of Ix*It
float sIyIt = 0.0;  // sum of Iy*It

for (int wy = -$window_radius; wy <= $window_radius; wy++) {
    for (int wx = -$window_radius; wx <= $window_radius; wx++) {
        float su = u + float(wx) * px;
        float sv = v + float(wy) * py;

        // Spatial gradients via central differences on current frame luminance
        float l_c = luma(sample(@current, su, sv));
        float l_r = luma(sample(@current, su + px, sv));
        float l_l = luma(sample(@current, su - px, sv));
        float l_u = luma(sample(@current, su, sv - py));
        float l_d = luma(sample(@current, su, sv + py));

        float gx = (l_r - l_l) * 0.5;
        float gy = (l_d - l_u) * 0.5;

        // Temporal gradient: difference between current and previous frame
        float l_prev = luma(sample(@previous, su, sv));
        float gt = l_c - l_prev;

        // Accumulate structure tensor
        sIx2  = sIx2  + gx * gx;
        sIy2  = sIy2  + gy * gy;
        sIxIy = sIxIy + gx * gy;
        sIxIt = sIxIt + gx * gt;
        sIyIt = sIyIt + gy * gt;
    }
}

// ── Solve 2x2 linear system ─────────────────────────────────────
// [sIx2   sIxIy] [vx]   [-sIxIt]
// [sIxIy  sIy2 ] [vy] = [-sIyIt]
//
// Cramer's rule: det = sIx2*sIy2 - sIxIy^2
float det = sIx2 * sIy2 - sIxIy * sIxIy;

// Corner response: det/trace ≈ harmonic mean of eigenvalues.
// High values indicate well-constrained flow (textured regions);
// low values indicate the aperture problem (edges) or flat regions.
float trace = sIx2 + sIy2;
float corner_response = sdiv(det, max(trace, 0.00001));
float confidence = smoothstep(0.0, $confidence_threshold, corner_response);

// Solve only if well-conditioned (det >> 0)
float inv_det = sdiv(1.0, det);
float vx = -(sIy2  * sIxIt - sIxIy * sIyIt) * inv_det;
float vy = -(sIx2  * sIyIt - sIxIy * sIxIt) * inv_det;

// Suppress flow in low-confidence regions
vx = vx * confidence;
vy = vy * confidence;

// Convert from UV displacement to pixel displacement
vx = vx * iw;
vy = vy * ih;

@OUT = vec4(vx, vy, confidence, 1.0);
